\section{Literature Review}

\subsection{Vehicle Rental Systems}

Traditional vehicle rental systems have long been characterized by manual booking processes, in-person identity verification, and centralized fleet management, each of which introduces operational inefficiencies and heightened susceptibility to fraud \cite{gupta2025double}. The academic response to these limitations has evolved along two principal trajectories. The first advocates for the digitization of identity verification workflows through KYC automation and document authentication pipelines, substantially reducing the overhead of manual screening while improving fraud detection accuracy. The second trajectory proposes the use of blockchain technology and Smart Contracts to establish immutable, self-executing rental agreements that eliminate reliance on centralized intermediaries \cite{madhusudan2019applying}. Protocols such as SePCAR further extend this framework by introducing Multi-Party Computation (MPC) for the generation of tamper-resistant access tokens, enabling secure vehicle handover even in offline conditions. While LuxeWay does not implement blockchain in its current version, these principles directly inform its emphasis on structured digital contracts, automated deposit handling, and transparent booking audit trails.

\subsection{Authentication and Security}

The security architecture of modern vehicle rental platforms is a well-established area of academic inquiry, driven by the growing volume of online transactions and the sensitivity of identity data exchanged during the rental lifecycle. Centralized authentication schemes are demonstrably vulnerable to credential stuffing, session hijacking, and unauthorized data access \cite{gupta2025double}. In response, the literature consistently advocates for stateless token-based authentication --- most commonly implemented via JSON Web Tokens (JWT) --- coupled with fine-grained role-based access control (RBAC) to constrain API surface exposure per user tier. LuxeWay operationalizes these recommendations directly: the backend enforces stateless JWT sessions with a 24-hour access token and a 7-day refresh token, while Spring Security's \texttt{HttpSecurity} configuration applies RBAC at the request-matcher level across four user roles (CUSTOMER, OWNER, ADMIN, SUPER\_ADMIN). OAuth2 social authentication via Google is additionally supported, with automatic email-based account linking to reduce onboarding friction without sacrificing identity assurance.

\subsection{Payment Systems}

The evolution of payment infrastructure in P2P rental platforms reflects a broader transition from cash-on-delivery and manual bank transfer to real-time, gateway-intermediated digital payments. In decentralized architectures, cryptocurrencies such as Ether enable direct peer-to-peer settlement, theoretically eliminating commission intermediaries entirely \cite{elhog2025decentralized}. Smart Contract escrow mechanisms further automate deposit release conditional on the fulfillment of predefined rental terms --- including on-time vehicle return and condition thresholds --- thereby reducing the need for manual dispute arbitration \cite{madhusudan2019applying}. LuxeWay adapts the escrow-and-release principle to the Vietnamese regulatory context by integrating \textbf{VNPay}'s payment gateway, which supports the full spectrum of domestic payment instruments (bank cards, MoMo, ZaloPay) alongside international Visa and Mastercard processing. The \textbf{LuxeWallet} internal balance further enables instant inter-party settlement without gateway round-trips. The pricing engine dynamically computes total booking cost across base rental, insurance tier selection, delivery fees, service charges, taxes, and coupon discounts, ensuring full pricing transparency prior to payment confirmation.

\subsection{Review and Rating Systems}

Reputation systems constitute the foundational trust layer of sharing-economy platforms, directly influencing booking conversion rates and platform retention \cite{hu2019sharing}. The inherent weakness of voluntary, text-based review systems --- susceptibility to fake reviews, reciprocal rating inflation, and selection bias among reviewers --- has been extensively documented in the literature. Proposed remedies range from algorithmic anomaly detection on review corpora to the integration of objective telematics signals sourced directly from vehicle OBD-II ports \cite{neifer2024trust}. Telematics-augmented reputation systems can evaluate driver behavior (speed compliance, braking patterns, mileage accuracy) independently of subjective reviewer opinion, substantially raising the bar for review manipulation. LuxeWay incorporates the conceptual insight of multi-signal trust evaluation by decomposing the post-rental review into five independently scored dimensions: vehicle cleanliness, listing accuracy, owner communication responsiveness, value for money, and delivery quality. Composite scores are aggregated and exposed on both vehicle listing pages and owner profiles, enabling prospective renters to make informed decisions grounded in a richer signal than a single aggregate star rating.
